% Fig 2.7 / spec Fig A03 -- cdf of a degenerate random variable
% Source: mathstatch2.tex lines 391-420
\documentclass[tikz,border=2pt]{standalone}
\usepackage{amsmath}
\usepackage{amssymb}
\usepackage{xcolor}
\usetikzlibrary{arrows.meta}
\newcommand{\sssig}{\scriptscriptstyle}

\definecolor{journalink}{HTML}{1F1C18}
\definecolor{journalmuted}{HTML}{8A7F72}
\definecolor{journalaccent}{HTML}{9C2D2D}

\begin{document}
\begin{tikzpicture}[
  x=1cm,
  y=1cm,
  every node/.style={font=\normalsize, journalink},
  axis/.style={journalink, line width=0.8pt, -{Stealth[length=4pt]}},
  tick/.style={journalink, line width=0.8pt},
  stepline/.style={journalink, line width=1.0pt},
  solidpt/.style={circle, fill=journalink, inner sep=0pt, minimum size=3.6pt},
  openpt/.style={circle, draw=journalink, line width=0.6pt, fill=white,
                 inner sep=0pt, minimum size=3.6pt}
]
  % ---- axes (x axis starts at -1 so the flat piece for x<0 is visible) ----
  \draw[axis] (-1,0) -- (3.5,0);
  \draw[axis] (0,0) -- (0,3.5);

  % ---- x ticks: 1.5 units per unit label; the tick at x=0 is lengthened
  %      to clear the open endpoint sitting there ----
  \draw[tick] (0,0.035) -- (0,-0.150);
  \foreach \x in {1.5,3}
    \draw[tick] (\x,0.035) -- (\x,-0.106);
  \foreach \x/\lab in {0/0, 1.5/1, 3/2}
    \node[anchor=north] at (\x,-0.25) {$\lab$};

  % ---- y ticks: only y=3, which corresponds to probability 1 ----
  \draw[tick] (0.035,3) -- (-0.106,3) node[anchor=east, inner sep=2pt] {$1$};

  % ---- step segments ----
  \draw[stepline] (-1,0) -- (0,0);
  \draw[stepline] (0,3)  -- (3,3);

  % ---- endpoints (drawn last) ----
  \node[openpt]  at (0,0) {};
  \node[solidpt] at (0,3) {};

  % ---- axis labels ----
  \node[anchor=west, inner sep=0pt] at (3.6,0) {$x$};
  \node[anchor=south, inner sep=0pt] at (0,3.6) {$F_{\sssig X}(x)$};
\end{tikzpicture}
\end{document}
